% Abstract / Kurzfassung

\newcommand{\hsmaabstracten}{%
PDF documents are the dominant format for sharing structured content across academia and professional domains, but they store pages as rendering instructions rather than semantic elements, making automated content extraction unreliable. This thesis designs, implements, and evaluates a multimodal PDF-to-Markdown conversion system that classifies each page individually into one of six content types (TEXT, IMAGE, FORMULA, TABLE, MIXED, or EMPTY) and routes each page to the most appropriate extraction strategy. Plain-text pages are processed by PyMuPDF without invoking a language model; visually complex pages are sent to a locally hosted vision-language model via LM Studio, eliminating any dependency on external cloud services. The system is evaluated on 90 pages from three bachelor thesis documents using a seven-metric framework covering text similarity, heading structure, list structure, table detection, code block accuracy, paragraph structure, and word overlap. Text-only extraction achieves the highest scores on five of seven metrics due to the predominantly prose-based nature of thesis documents. The adaptive strategy provides a measurable improvement in table structure accuracy (0.956 versus 0.922) and reduces total token consumption by 71\% compared to image-only processing. A two-stage post-processing pipeline removes LLM output artefacts and reconstructs document structure, verified by 147 unit tests. The results show that per-page strategy selection is more efficient than any uniform fixed strategy, though text-based extraction remains the strongest baseline for text-dominant documents.%
}

\newcommand{\hsmaabstractde}{%
PDF-Dokumente sind das dominierende Format für den Austausch strukturierter Inhalte in Wissenschaft und Praxis, speichern Seiten jedoch als Rendering-Anweisungen statt als semantische Elemente, was eine automatisierte Inhaltsextraktion erschwert. Diese Bachelorarbeit entwirft, implementiert und evaluiert ein multimodales System zur Konvertierung von PDF-Dokumenten in strukturiertes Markdown, das jede Seite einzeln in einen von sechs Inhaltstypen klassifiziert (TEXT, IMAGE, FORMULA, TABLE, MIXED oder EMPTY) und an die jeweils geeignetste Extraktionsstrategie weiterleitet. Reine Textseiten werden von PyMuPDF ohne LLM-Aufruf verarbeitet; visuell komplexe Seiten werden an ein lokal gehostetes Vision-Language-Modell über LM Studio übergeben, ohne Abhängigkeit von externen Cloud-Diensten. Das System wird auf 90 Seiten aus drei Bachelorarbeiten mit einem Sieben-Metrik-Rahmen evaluiert, der Textähnlichkeit, Überschriftenstruktur, Listenstruktur, Tabellenerkennung, Codeblock-Genauigkeit, Absatzstruktur und Wortüberschneidung abdeckt. Die reine Textextraktion erzielt auf fünf von sieben Metriken die höchsten Werte. Die adaptive Strategie liefert messbare Verbesserungen bei der Tabellenstruktur (0,956 gegenüber 0,922) und reduziert den gesamten Token-Verbrauch um 71\,\% gegenüber der reinen Bildextraktion. Eine zweistufige Nachverarbeitungspipeline entfernt LLM-Ausgabe-Artefakte und rekonstruiert die Dokumentstruktur, überprüft durch 147 Unit-Tests. Die Ergebnisse zeigen, dass die seitenweise Strategiewahl effizienter ist als jede einheitliche Fixstrategie, wobei die textbasierte Extraktion die stärkste Baseline für textdominante Dokumente bleibt.%
}
